\section{A graph of large distance classes}\label{sec:extremal}

Suppose that the conclusion of Theorem~\ref{thm:main} fails. We first
extract a positive-density directed graph whose edges belong to large
distance classes. The extraction also controls how much of a class can
lie in a small subset of vertices. These two properties will allow us to
compare distance classes with the uniform distribution on the whole set.

Put
\[
  \Theta(t)=\limsup_{n\to\infty}F_n(t),\qquad
  s_* = \sup\{t>0:\Theta(t)>0\}.
\]
Each function $F_n$, and hence $\Theta$, is nonincreasing. Moreover,
$F_n(t)=0$ for $t\ge1$, since a distance class has at most $n-1$ points.
Our supposition therefore gives $0<s_*\le1$. Fix
\begin{equation}\label{eq:ext-parameters}
  0<c<s_*/3,\qquad
  \max\{c,(1-c)s_*\}<s<s_*,\qquad c<\gamma<s,
\end{equation}
choosing $s$ to be a continuity point of $\Theta$. Such a choice is
possible because the indicated interval is nonempty and a monotone
function has at most countably many discontinuities. Set
$\theta=\Theta(s)>0$; positivity follows from $s<s_*$ and monotonicity.
These parameters remain fixed throughout the proof. All asymptotic
estimates below refer to the sequence of configurations constructed in
Proposition~\ref{prop:graph}, or to a subsequence of it. Their constants
may depend on these fixed parameters, but not on the configurations or
their coordinates.

We record the uniformity consequence of continuity that will be used
twice in the extraction.
\begin{lemma}\label{lem:ext-tail}
If $r_N\to s$ and $M_N\to\infty$, then
\[
  \sup_{m\ge M_N}(F_m(r_N)-\theta)_+\longrightarrow0,
\]
where the supremum is over integer $m\ge2$ and $(u)_+=\max\{u,0\}$.
\end{lemma}
\begin{proof}
Given $\epsilon>0$, continuity permits a fixed $\delta>0$ such that
$\Theta(s-\delta)<\theta+\epsilon/2$. For sufficiently large $m$,
$F_m(s-\delta)\le\Theta(s-\delta)+\epsilon/2$, by the definition of
limsup. For sufficiently large $N$, we have $r_N\ge s-\delta$ and all
$m\ge M_N$ are in this range. Monotonicity now bounds every
$F_m(r_N)$ by $\theta+\epsilon$.
\end{proof}

Here are the probability laws associated with a nonempty directed graph
$E\subset\{(x,y)\in P^2:x\ne y\}$ on an $n$-point set $P$. Partition
its edges into \emph{fibers}: edges lie in the same fiber when they have
the same source and the same squared distance. Write $x_e$ for the
source of a fiber $e$, $F_e\subset P\setminus\{x_e\}$ for its target
set, and $k_e=|F_e|$. Let $d^-(y)$ denote the number of incoming edges at
$y$.
\begin{center}
\renewcommand{\arraystretch}{1.2}
\begin{tabular}{@{}p{0.18\textwidth}p{0.75\textwidth}@{}}
\hline
Law & Definition \\
\hline
$\alpha$ & Uniform probability measure on $P$. \\
$\pi$ & Uniform probability measure on the $n(n-1)$ distinct ordered
  pairs in $P$. \\
$q$ & Uniform probability measure on $E$, with source and target
  marginals $q_X$ and $q_Y$. \\
$\mu_e$ & Uniform probability measure on $F_e$. \\
$\rho_y$ & Uniform probability measure on the incoming neighbors of
  $y$, defined when $d^-(y)>0$. \\
\hline
\end{tabular}
\end{center}
To sample $q$, one can first choose $e$ with probability $k_e/|E|$ and
then choose $y$ according to $\mu_e$, setting $x=x_e$. We write
$\E_e$ for expectation with these fiber weights. Equivalently, one can
choose $y$ according to $q_Y$ and then $x$ according to $\rho_y$.
All these laws are attached to the finite graph; they do not depend on
any of the absolute values introduced later.

\begin{proposition}\label{prop:graph}
Under the supposition above, there is a sequence of finite sets
$P\subset\R^2$ of real algebraic points, with $n=|P|\to\infty$, and
nonempty directed graphs $E$ on $P$, with the following properties.
Their densities $d_n=|E|/[n(n-1)]$ tend to $\theta$, and every fiber has
$k_e\ge n^{s-o(1)}$, uniformly in $e$. For a fixed constant $B_0$,
\begin{equation}\label{eq:ext-domination}
  q\le B_0\pi,\qquad q_X,q_Y\le B_0\alpha,\qquad
  \frac{q_X+q_Y}{2}\ge(1-o(1))\alpha
\end{equation}
pointwise, and
\begin{equation}\label{eq:ext-small-degree}
  q_Y\{y:d^-(y)<n^\gamma\}=O(n^{\gamma-1})=o(1).
\end{equation}
Finally, put $\lambda_n=1/\log n$. For every subset $C\subset P$, with
$m=|C|$ and $p_e=\mu_e(C)$,
\begin{equation}\label{eq:subset}
 \E_e\,\1_C(x_e)p_e\1_{\{p_e\ge\lambda_n\}}
 \le
 \begin{cases}
   (1+o(1))\pi(C^2),&m>n^{1-c},\\
   o(1)\pi(C^2),&m\le n^{1-c}.
 \end{cases}
\end{equation}
All the errors are uniform in vertices, fibers, and subsets wherever
these occur.
\end{proposition}

\begin{proof}
We construct the graph, prove that deleting low-degree vertices cannot
remove half the configuration, and then prove the subset estimate.

\paragraph{Preserving the exact distance pattern.}
For fixed $N$, the numerator defining $F_N(s)$ takes values in the
finite set $\{0,1,\ldots,N(N-1)\}$. Its maximum is therefore attained,
without any compactness assertion about configurations. Choose a
sequence $N\to\infty$ on which $F_N(s)\to\theta$ and an attaining
configuration for each $N$.

Each such configuration can be replaced by one with real algebraic
coordinates without changing any distance class. To see this, label its
points by $1,\ldots,N$ and use coordinate variables $(u_i,b_i)$. Write
$D_{ij}=(u_i-u_j)^2+(b_i-b_j)^2$. Require $D_{ij}>0$ for $i<j$,
and prescribe, for every two unordered pairs, the equality
$D_{ij}=D_{kl}$ or the disequality $D_{ij}\ne D_{kl}$ that holds in
the original configuration. These are finitely many polynomial
conditions over $\Q$; a disequality can be expressed as
$(D_{ij}-D_{kl})^2>0$. Their existential sentence holds over $\R$ and
therefore over the real algebraic numbers, by transfer for real closed
fields~\cite[Theorems~3.1 and~4.1]{RCFTransfer}. The resulting points are
distinct and retain the exact distance equality pattern. In particular every class size,
and hence the value attaining $F_N(s)$, is unchanged. The real threshold
$N^s$ need not enter the transferred system: its comparison with each
preserved integer class size is unchanged. No bound on the
degrees or heights of these coordinates is imposed.

\paragraph{Deleting vertices and small fibers.}
On this algebraic configuration, initially keep every edge counted by
$F_N(s)$. Each initial fiber has at least $N^s$ edges. Define
\[
 K_N=\frac{N^s}{\log N},\qquad
 f_N=\frac{\log K_N}{\log N}
     =s-\frac{\log\log N}{\log N},\qquad
 m_0=\lceil N/2\rceil.
\]
For all sufficiently large $N$, $f_N>0$. By
Lemma~\ref{lem:ext-tail},
\[
 b_N:=\max_{m_0\le m\le N}(F_m(f_N)-\theta)_+\longrightarrow0.
\]
Put
\[
 t_N=b_N+|F_N(s)-\theta|+\frac1{\log N},\qquad
 \eta_N=\sqrt{t_N}.
\]
Thus $\eta_N\to0$, and eventually $\theta-\eta_N>0$.

Whenever a remaining nonempty fiber has fewer than $K_N$ edges, discard
it. Also, if the current vertex count is $m$ and a vertex has the sum
of its outgoing and incoming degrees less than
$2(\theta-\eta_N)(m-1)$, delete that vertex and all its incident edges.
Clean out small fibers after each vertex deletion, and continue until
all degree lower bounds hold.

There are at most $N(N-1)/N^s$ initial fibers. Restricting vertices
never splits a source-distance fiber, and each initial fiber can be
discarded only once, at a cost below $K_N$. Thus all fiber cleanups
together cost at most $N(N-1)/\log N$ edges. Suppose that the process
reaches $m_0$ vertices. Put $A_N=N(N-1)$ and
$B_N=m_0(m_0-1)$. The total cost of vertex deletions is less than
\[
 2(\theta-\eta_N)\sum_{m=m_0+1}^{N}(m-1)
 = (\theta-\eta_N)(A_N-B_N).
\]
After cleanup, the number of remaining edges is therefore at least
\[
 F_N(s)A_N-(\theta-\eta_N)(A_N-B_N)-\frac{A_N}{\log N}.
\]
Its excess over $(\theta+b_N)B_N$ is at least
\[
 \eta_N(A_N-B_N)-t_NA_N
 =\eta_NA_N\left(1-\frac{B_N}{A_N}-\eta_N\right)>0
\]
for large $N$, because $B_N/A_N\to1/4$. On the other hand, every
remaining edge belongs to a fiber of size at least
$K_N=N^{f_N}\ge m_0^{f_N}$. Its density is therefore bounded above by
$F_{m_0}(f_N)\le\theta+b_N$, a contradiction.

The process consequently stops on a set $P$ of size $n>m_0$, with a
graph $E$ whose nonempty fibers have size at least $K_N$. Summing the
degree lower bounds, and applying the definition of $b_N$ to the
remaining graph, gives
\begin{equation}\label{eq:ext-density}
  \theta-\eta_N\le d_n\le F_n(f_N)\le\theta+b_N.
\end{equation}
Here the upper bound applies because $K_N\ge n^{f_N}$. Thus
$d_n\to\theta$ and $k_e\ge K_N\ge n^{s-o(1)}$ uniformly in $e$.

\paragraph{The probability bounds.}
For every edge, $q(x,y)/\pi(x,y)=1/d_n$, so $B_0=2/\theta$ works
in the first inequality of \eqref{eq:ext-domination} for large $n$.
Summing over one coordinate gives the upper marginal bounds. At every
vertex, the degree lower bound gives
\[
 \frac{q_X(x)+q_Y(x)}2
 \ge\frac{\theta-\eta_N}{d_n}\alpha(x)
 \ge\frac{\theta-\eta_N}{\theta+b_N}\alpha(x)
 =(1-o(1))\alpha(x).
\]
Only the sum of the two marginals is bounded below; separate lower
bounds are not needed. The targets of small incoming degree satisfy
\[
 q_Y\{y:d^-(y)<n^\gamma\}
 \le\frac{n\,n^\gamma}{|E|}=O(n^{\gamma-1}),
\]
which proves \eqref{eq:ext-small-degree} since $\gamma<s<1$.

\paragraph{Large fibers inside subsets.}
It remains to prove \eqref{eq:subset}. Define
\[
 T_n=\lambda_nK_N=\frac{N^s}{\log N\log n}.
\]
Since $N/2<n\le N$, there is a sequence $\delta_n\to0$ such that
$T_n\ge n^{s-\delta_n}$; in particular, $T_n\to\infty$.
Multiplying the left side of \eqref{eq:subset} by $|E|$ counts exactly
the edges $(x_e,y)\in C^2$ for which
$|F_e\cap C|\ge\lambda_n k_e$. Every counted pair therefore satisfies
\[
 k_C(x_e,y)\ge |F_e\cap C|\ge T_n.
\]
A nonzero contribution requires $m-1\ge T_n$, because the source is
in $C$ and is excluded from its own distance class. Thus all bounded
subsets contribute zero for sufficiently large $n$; the cases $m=0,1$
give zero at every $n$.

For $m>n^{1-c}$, the positive exponent $s-\delta_n$ gives
$T_n\ge n^{s-\delta_n}\ge m^{s-\delta_n}$. By
Lemma~\ref{lem:ext-tail}, $F_m(s-\delta_n)\le\theta+o(1)$ uniformly
in all these sizes. Dividing the resulting pair count by $|E|$ yields
\[
 \E_e\,\1_C(x_e)p_e\1_{\{p_e\ge\lambda_n\}}
 \le\frac{\theta+o(1)}{d_n}
       \frac{m(m-1)}{n(n-1)}
 =(1+o(1))\pi(C^2).
\]

For $m\le n^{1-c}$, choose a fixed $s'$ with
$s_*<s'<s/(1-c)$, possible by \eqref{eq:ext-parameters}.
For sufficiently large $n$,
\[
 T_n\ge n^{s-\delta_n}\ge n^{(1-c)s'}\ge m^{s'}.
\]
The definition of $s_*$ gives $\Theta(s')=0$, hence $F_m(s')\to0$
as $m\to\infty$. Every contributing subset has $m\ge T_n+1$, so
\[
 \sup_{m\ge T_n+1}F_m(s')\longrightarrow0,
\]
where again only integer sizes are considered. Dividing this pair
bound by $|E|$ gives $o(1)\pi(C^2)$ uniformly over all contributing
small subsets; the other subsets contribute zero. This also covers
$s'>1$, when the corresponding pair count is identically zero.
\end{proof}

We henceforth work with the sequence in Proposition~\ref{prop:graph}.
In particular, every $k_e\ge n^\gamma$ for sufficiently large $n$.
The graph and its sampling laws are now fixed. The next section builds
partitions adapted to the two factors of squared Euclidean distance;
Equation~\eqref{eq:subset} will later control their small cells.
